%%----------------------------------------------------------------------------
%%----------------------------------------------------------------------------
\section{Proyecto final: Proyecto libre (2026)}
\label{practica-final-2026-05-ltaw}

\textbf{Repositorio plantilla.}

\url{https://gitlab.eif.urjc.es/cursosweb/2025-2026-ltaw/final}

[ \textbf{Nota importante:} Enunciado aún en proceso de elaboración, puede sufrir algunos cambios antes de su versión final. Por favor, avisa si detectas algún error o inconsistencia. Al final de este enunciado puede verse una lista de preguntas frecuentes (apartado~\ref{sec:practica-2026-05-ltaw:preguntas}). ]

Este es el enunciado del proyecto final del curso 2025-2026, tanto para la convocatoria ordinaria como para la extraordinaria.

A diferencia de ediciones anteriores, el proyecto final de este curso es de temática \textbf{libre}: cada alumno propondrá y desarrollará su propia aplicación web, aplicando los conocimientos adquiridos a lo largo de la asignatura. El objetivo es demostrar dominio real de las tecnologías del curso en un proyecto genuinamente propio.

La práctica es de carácter \textbf{individual}.

A continuación se describen los requisitos generales, el proceso de seguimiento con los profesores, los criterios de evaluación y algunas ideas de proyecto a modo de orientación.

%%----------------------------------------------------------------------------
\subsection{Requisitos generales}

El proyecto deberá:

\begin{itemize}

\item Construirse como un proyecto Django/Python3, con una o varias aplicaciones (\emph{apps}) Django que implementen la funcionalidad requerida.

\item Usar SQLite3 para el almacenamiento persistente, con tablas definidas mediante modelos Django.

\item Incluir acceso a todos los datos relevantes a través del ``Admin Site'' de Django.

\item Utilizar plantillas Django con jerarquía de herencia para generar las páginas HTML.

\item Cubrir entre el 70\,\% y el 80\,\% de los contenidos vistos en la asignatura (véase el apartado de evaluación para más detalles).

\item Integrar al menos una \textbf{API externa real} (de datos abiertos, de inteligencia artificial generativa, de servicios de terceros, etc.).

\item Seguir las convenciones de marcado HTML, CSS y Bootstrap (u otra biblioteca de estilo) usadas en la asignatura: elementos con atributo \texttt{id}, clases CSS compartidas y diseño \emph{responsive} que funcione correctamente en escritorio y móvil.

\item Incluir cabecera, menú de navegación y pie de página comunes a todas las páginas de la aplicación.

\end{itemize}

Además, se valorará positivamente el uso de:

\begin{itemize}
\item HTMX para actualización dinámica de contenidos.
\item Tests extremo a extremo y/o unitarios con buena cobertura.
\item Recursos en formato JSON útiles para la aplicación.
\item Autenticación de usuarios mediante el sistema de sesiones de Django.
\end{itemize}

%%----------------------------------------------------------------------------
\subsection{Proceso de trabajo y seguimiento con los profesores}
\label{sec:practica-2026-05-ltaw:fases}

Dado el carácter libre del proyecto, el seguimiento con el profesorado es \textbf{obligatorio} y forma parte de la evaluación. El proceso se organiza en las siguientes fases, que deberán quedar documentadas en el fork del repositorio de plantilla:

\begin{description}

\item[Fase 0 -- Propuesta inicial.] Antes de comenzar el desarrollo, el alumno deberá presentar a los profesores una propuesta escrita del proyecto. La propuesta se entregará creando un fichero \texttt{PROPUESTA.md} en el repositorio del fork y notificando a los profesores. Deberá incluir al menos:
  \begin{itemize}
    \item Descripción general de la aplicación y su utilidad.
    \item Tecnologías y APIs externas que se van a utilizar.
    \item Esquema preliminar del modelo de datos (tablas principales).
  \end{itemize}
  Los profesores darán su visto bueno o propondrán ajustes antes de que el alumno comience el desarrollo. \textbf{No se puede empezar a desarrollar sin aprobación de la Fase 0.}

\item[Fase 1 -- Prototipo funcional.] El alumno presentará a los profesores un prototipo con la estructura básica de la aplicación: modelo de datos implementado, al menos dos o tres recursos web funcionando y la integración inicial con la API externa. Esta revisión se documentará en el fichero \texttt{FASES.md} del repositorio, indicando lo completado, lo pendiente y los cambios respecto a la propuesta inicial.

\item[Fase 2 -- Versión avanzada.] Se presentará una versión con la funcionalidad obligatoria completa y, si se ha implementado, la parte optativa. Se actualizará igualmente el fichero \texttt{FASES.md}.

\item[Entrega final.] La entrega incluirá el repositorio completo, los vídeos de demostración y el \texttt{README.md} con todos los datos requeridos (ver apartado de entrega).

\end{description}

Las fechas de revisión de las fases intermedias se acordarán con los profesores al aprobar la propuesta. El incumplimiento de las fases intermedias puede penalizar la nota del proyecto.

Es importante que cada fase esté \textbf{definida con la mayor precisión posible}, especificando avances concretos y medibles: qué funcionalidad estará terminada, qué recursos estarán implementados y qué se podrá probar en ese momento. Las descripciones vagas (``avanzado'', ``casi listo'', ``en progreso'') no se considerarán válidas como cierre de fase. Los profesores valorarán la calidad de la planificación y el rigor con el que se documentan los avances en \texttt{FASES.md}.

%%----------------------------------------------------------------------------
\subsection{Evaluación}
\label{sec:practica-2026-05-ltaw:evaluacion}

Según la guía docente, la nota del proyecto final se compone de:

\begin{description}

\item[Parte obligatoria (0--2 puntos).] Se evaluará que el proyecto:
  \begin{itemize}
    \item Cubre entre el 70\,\% y el 80\,\% de los contenidos de la asignatura (con desglose en el README).
    \item Incluye al menos una integración real con una API externa.
    \item Funciona correctamente en el entorno del laboratorio y en el despliegue.
    \item Ha seguido el proceso de fases descrito anteriormente, con la documentación correspondiente en \texttt{PROPUESTA.md} y \texttt{FASES.md}.
    \item Incluye un README completo con todos los datos requeridos.
  \end{itemize}

\item[Parte optativa (0--3 puntos).] Queda a criterio del profesor, en función de la calidad, originalidad e impacto de lo desarrollado más allá de lo mínimo obligatorio. El alumno deberá indicar explícitamente en el README qué considera parte optativa y por qué. Algunos ejemplos de lo que puede considerarse optativo:
  \begin{itemize}
    \item Uso de HTMX para funcionalidades dinámicas.
    \item Tests con buena cobertura (unitarios y/o extremo a extremo).
    \item Integración con múltiples APIs externas o uso avanzado de una API.
    \item Calidad visual destacable y experiencia de usuario cuidada.
    \item Participación en los II Premios a la Reutilización de Datos Abiertos 2026 (ver apartado~\ref{sec:practica-2026-05-ltaw:apis}).
    \item Cualquier otra funcionalidad que el alumno justifique como relevante.
  \end{itemize}

\end{description}

\textbf{Nota total del proyecto: 0--4 puntos (mínimo 2 para superar el proyecto).} La nota total se calcula a partir de ambas partes, pero está limitada a 4 puntos. Obtener 0 en la parte obligatoria implica no superar el proyecto final.

%%----------------------------------------------------------------------------
\subsection{Ideas de proyectos}
\label{sec:practica-2026-05-ltaw:ideas}

A continuación se proponen tres ideas concretas de proyecto, a modo de orientación. Son sugerencias: el alumno puede proponer cualquier otra idea siempre que cumpla los requisitos generales y sea aprobada por los profesores en la Fase 0.

\begin{description}

\item[Corrector automático de prácticas.] Aplicación web que utilice la API de GitLab de la EIF\footnote{API de GitLab: \url{https://docs.gitlab.com/ce/api/}} para descargarse automáticamente los forks de una plantilla de práctica y emplee la API de un LLM (por ejemplo, NVIDIA Build) para analizar y corregir el código de cada fork. La aplicación mostraría un informe de corrección por alumno con comentarios generados por el modelo. Resultaría especialmente útil para el propio profesorado de la asignatura.

\item[Generador de CVN.] Aplicación web que, a partir de datos introducidos por el usuario o recuperados de APIs públicas, genere currículums vitae en formato CVN (Currículum Vítae Normalizado del Ministerio de Ciencia\footnote{CVN: \url{https://cvn.fecyt.es/}}), el estándar utilizado en convocatorias de proyectos públicos en España.

\item[Explorador de datos abiertos.] Aplicación web que permita explorar, visualizar y analizar conjuntos de datos abiertos publicados por la Administración Pública. Como referencia se propone el portal de datos abiertos de la Comunidad de Madrid\footnote{Datos abiertos de la Comunidad de Madrid: \url{https://datos.madrid.es/}} y el del Gobierno de España\footnote{Datos.gob.es: \url{https://datos.gob.es/}}, aunque se puede usar cualquier portal equivalente. La aplicación podría incluir búsqueda, filtrado, representaciones gráficas y exportación de resultados.

  Este proyecto es especialmente interesante porque podría presentarse a los \textbf{II Premios a la Reutilización de Datos Abiertos (2026)}\footnote{\url{https://datos.madrid.es/pages/premios-de-reutilizacion-2026}} del Ayuntamiento de Madrid. La participación en estos premios se valorará como parte optativa.

\end{description}

%%----------------------------------------------------------------------------
\subsection{APIs e integración con servicios externos}
\label{sec:practica-2026-05-ltaw:apis}

El proyecto debe integrar al menos una API externa. A continuación se proporcionan algunos recursos de referencia.

\subsubsection*{Datos abiertos}

\begin{itemize}
  \item Portal de datos abiertos del Ayuntamiento de Madrid: \url{https://datos.madrid.es/}
  \item Portal de datos abiertos de la Comunidad de Madrid: \url{https://www.comunidad.madrid/gobierno/datos-abiertos}
  \item Portal de datos abiertos del Gobierno de España: \url{https://datos.gob.es/}
\end{itemize}

\subsubsection*{Modelos de lenguaje (LLM)}

Si el proyecto requiere capacidades de inteligencia artificial generativa, se propone utilizar NVIDIA Build\footnote{\url{https://build.nvidia.com/}} como proveedor recomendado, ya que ofrece modelos gratuitos accesibles mediante API. También se permite usar cualquier otro proveedor con acceso por API (por ejemplo, un modelo local con LM Studio, OpenRouter u otros).

Requisitos para cualquier integración con LLM:

\begin{itemize}
  \item El token o clave de API \textbf{no} debe subirse al repositorio.
  \item Debe gestionarse mediante variable de entorno (por ejemplo, \verb|NVIDIA_API_KEY|) o mediante fichero \verb|.env| incluido en \verb|.gitignore|.
\end{itemize}

En la web de NVIDIA Build encontraréis el detalle actualizado de uso de la API y los parámetros disponibles en cada modelo.

Ejemplo recomendado (usando \verb|requests|):

\begin{verbatim}
import os
import requests
import json


def main():
  url = "https://integrate.api.nvidia.com/v1/chat/completions"

  headers = {
    "Authorization": f"Bearer {os.getenv('NVIDIA_API_KEY')}",
    "Content-Type": "application/json",
  }

  payload = {
    "model": "google/gemma-2-2b-it",
    "messages": [
      {
        "role": "user",
        "content": "Write a Go program that generates the Fibonacci sequence and prints the first 10 numbers."
      }
    ],
    "temperature": 0.2,
    "top_p": 0.7,
    "max_tokens": 1024,
    "stream": True,
  }

  response = requests.post(url, headers=headers, json=payload, stream=True)

  for line in response.iter_lines():
    if line:
      decoded = line.decode("utf-8")

      # El streaming viene en formato SSE: "data: {...}"
      if decoded.startswith("data: "):
        data = decoded[len("data: "):]

        if data == "[DONE]":
          break

        try:
          chunk = json.loads(data)
          delta = chunk["choices"][0]["delta"]

          if "content" in delta:
            print(delta["content"], end="")

        except json.JSONDecodeError:
          continue


if __name__ == "__main__":
  main()
\end{verbatim}

Alternativamente, en la propia web de NVIDIA suelen mostrar ejemplos con el cliente de OpenAI. También podéis usar ese enfoque si os resulta más cómodo:

\begin{verbatim}
import os
from openai import OpenAI


def main() -> None:
    client = OpenAI(
        base_url="https://integrate.api.nvidia.com/v1",
        api_key=os.getenv("NVIDIA_API_KEY"),
    )

    completion = client.chat.completions.create(
        model="google/gemma-2-2b-it",
        messages=[{"role": "user", "content": "Hola, ¿cómo estás?"}],
        temperature=0.2,
        top_p=0.7,
        max_tokens=1024,
        stream=True,
    )

    for chunk in completion:
        if chunk.choices and chunk.choices[0].delta.content is not None:
            print(chunk.choices[0].delta.content, end="")


if __name__ == "__main__":
    main()
\end{verbatim}

En ese caso, recordad incluir también \verb|openai| en \verb|requirements.txt|.

%%----------------------------------------------------------------------------
\subsection{Despliegue}
\label{sec:practica-2026-05-ltaw:despliegue}

La práctica deberá estar desplegada en algún sitio de Internet accesible. Deberá mantenerse activa al menos desde el día de entrega de la práctica hasta el cierre de actas.

Para el despliegue se puede utilizar Python Anywhere\footnote{Python Anywhere: \url{https://pythonanywhere.com}}, que proporciona un plan gratuito suficiente para este proyecto. También se puede desplegar en un ordenador dedicado (por ejemplo, una Raspberry Pi accesible desde Internet) o en servicios como Google Computing Engine\footnote{GCP Engine Free: \url{https://cloud.google.com/free/}}. Este tipo de despliegues, al requerir más configuración propia, estarán algo más valorados.

En el caso de que la práctica se despliegue en Python Anywhere, hay que tener en cuenta que sus máquinas virtuales tienen cortado el acceso a todos los sitios de Internet salvo los que están en una ``lista blanca''\footnote{Lista blanca de Python Anywhere: \url{https://www.pythonanywhere.com/whitelist/}} (\emph{whitelist}). Si vuestro proyecto se conecta a APIs externas, verificad que dichas APIs están en esa lista. Si no lo están, aseguraos de que la base de datos del despliegue ya incluya los datos necesarios.

Tened en cuenta que si usáis otras plataformas para el despliegue, puede que os encontréis problemas similares. En cualquier caso, los profesores probaremos la práctica en el despliegue, así que todos los recursos implementados deben funcionar correctamente si no hay restricciones de conexión.

Tenéis más detalles sobre cómo se hace un despliegue de una aplicación Django en Python Anywhere en el vídeo ``Django: Despliegue en Python Anywhere''\footnote{\url{https://www.youtube.com/watch?v=hlZPC5L2Itc}}, que explica cómo desplegar allí la aplicación \texttt{django-youtube-4}\footnote{\url{https://github.com/CursosWeb/Code/tree/master/Python-Django/django-youtube-4}} (ver también el fichero \texttt{README.md} de esa aplicación para más detalles).

%%----------------------------------------------------------------------------
\subsection{Entrega de la práctica}
\label{sec:practica-2026-05-ltaw:entrega}

\begin{itemize}
  \item \textbf{Fecha límite de entrega (convocatoria ordinaria):} día anterior al del examen de la asignatura a las 23:55 (hora española peninsular)
  \item \textbf{Fecha límite de entrega (convocatoria extraordinaria):} día anterior al del examen de la asignatura, a las 23:55 (hora española peninsular)
  \item \textbf{Realización de entrevistas:} junto con la revisión de la asignatura.
\end{itemize}

La entrega de la práctica consiste en:

\begin{enumerate}

  \item {\bf Subir tu práctica a un repositorio en el GitLab de la Escuela}. El repositorio contendrá todos los ficheros necesarios para que funcione la aplicación (ver detalle más abajo). Es imprescindible que el alumno haya realizado un fork del repositorio de plantilla indicado al principio de este enunciado.

El repositorio deberá incluir, además del código y la base de datos, los ficheros \texttt{PROPUESTA.md} y \texttt{FASES.md} con el historial del seguimiento del proyecto.

Recordad que es importante ir haciendo commits con frecuencia y hacer push antes de la entrega. Se recomienda mantener el repositorio como privado hasta el momento de la entrega.

\item Para considerar la práctica como entregada es fundamental que en el directorio raíz del repositorio haya un fichero \texttt{README.md} cuya primera línea sea:

\begin{verbatim}
# Entrega convocatoria XXX
\end{verbatim}

Siendo \texttt{XXX} ``mayo'' si se entrega en la convocatoria ordinaria de mayo, y ``junio'' si se entrega en la convocatoria extraordinaria de junio.

\textbf{Atención:} No se considerará entregada la práctica si no ve esta primera línea en el formato adecuado.

 \item {\bf Entregar un vídeo de demostración de la parte obligatoria, y otro vídeo de demostración de la parte opcional}, si se ha realizado parte optativa. Los vídeos serán de una {\bf duración máxima de 3 minutos} (cada uno), y consistirán en una captura de pantalla de un navegador web utilizando la aplicación, y mostrando lo mejor posible la funcionalidad correspondiente (básica u opcional). Siempre que sea posible, el alumno comentará en el audio del vídeo lo que vaya ocurriendo en la captura. Los vídeos se colocarán en algún servicio de subida de vídeos en Internet (por ejemplo, Vimeo, Twitch, o YouTube). Los vídeos de más de tres minutos tendrán penalización.

Hay muchas herramientas que permiten realizar la captura de pantalla. Por ejemplo, en GNU/Linux puede usarse Gtk-RecordMyDesktop o Istanbul (ambas disponibles en Ubuntu). OBS Studio\footnote{OBS Studio: \url{https://obsproject.com/}} está disponible para varias plataformas (Linux, Windows, MacOS). Es importante que la captura sea realizada de forma que se distinga razonablemente lo que se grabe en el vídeo.

En caso de que convenga editar el vídeo resultante (por ejemplo, para eliminar tiempos de espera) puede usarse un editor de vídeo, pero siempre deberá ser indicado que se ha hecho tal cosa con un comentario en el audio, o un texto en el vídeo. Hay muchas herramientas que permiten realizar esta edición. Por ejemplo, en GNU/Linux puede usarse OpenShot o PiTiVi.

\end{enumerate}

Sobre el contenido del repositorio:
\begin{itemize}
  \item Se han de entregar los siguientes ficheros:

\begin{itemize}
\item El repositorio en la instancia GitLab de la EIF deberá tener, en su directorio raíz, un proyecto Django completo y listo para funcionar en el entorno del laboratorio, incluyendo la base de datos. Deberá poder ejecutarse directamente con \verb|python3 manage.py runserver| desde un entorno virtual en el que esté instalado Django. También ejecutará los tests con \verb|python3 manage.py test|, desde el mismo entorno virtual.

\item La base de datos habrá de tener datos suficientes como para poder evaluar la aplicación, incluyendo datos obtenidos de las APIs externas utilizadas.

\item Un fichero \verb|requirements.txt|, con un nombre de paquete Python por línea, para indicar cualquier biblioteca Python necesaria para que funcione la aplicación, incluyendo Django. En el formato compatible con \verb|pip install -r requirements.txt|.

\item Los ficheros \texttt{PROPUESTA.md} y \texttt{FASES.md} con la documentación del proceso de desarrollo.

\item Cualquier otro fichero auxiliar que pueda hacer falta para que funcione la práctica.
\end{itemize}

\item Se incluirán en el fichero README.md los siguientes datos (atención a lo indicado sobre la primera línea de este fichero):

\begin{itemize}
  \item Nombre y titulación.
  \item Nombre de cuenta en laboratorios Linux de la EIF.
  \item Nombre de cuenta de la URJC.
  \item URL del vídeo demostración de la funcionalidad básica.
  \item URL del vídeo demostración de la funcionalidad optativa, si se ha realizado funcionalidad optativa.
  \item URL de la aplicación desplegada.
  \item Al menos un usuario y contraseña que permita autenticarse para poder usar la aplicación.
  \item Cuenta del superusuario para entrar en el Admin Site.
  \item Listado de recursos implementados y métodos HTTP permitidos en cada uno.
  \item Resumen de lo implementado en la parte obligatoria, indicando qué contenidos de la asignatura se han aplicado y en qué medida se estima que se cubre el 70--80\,\% del temario.
  \item Lista de funcionalidades optativas que se hayan implementado, con breve descripción de cada una y justificación de por qué se consideran optativas.
\end{itemize}

Estos datos se escribirán siguiendo estrictamente el siguiente formato (si se entrega en junio, el título será ``Entrega práctica junio'':

\begin{verbatim}
# Entrega práctica mayo

## Datos

* Nombre:
* Titulación:
* Cuenta en laboratorios:
* Cuenta URJC:
* Vídeo básico (URL):
* Vídeo parte opcional (URL):
* Despliegue (URL):
* Usuarios y contraseñas:
* Cuenta Admin Site: usuario/contraseña

## Recursos y métodos HTTP

* Recurso: /ruta/ejemplo
* Métodos permitidos: GET, POST
* Recurso: /ruta/otra
* Métodos permitidos: GET

## Resumen parte obligatoria

(Descripción de lo implementado e indicación de qué contenidos
del curso se han aplicado y en qué medida se cubre el 70-80%)

## Lista partes opcionales

* Nombre parte: descripción y justificación
* ...
\end{verbatim}

Asegúrate de que el fichero esté adecuadamente formateado en Markdown.
\end{itemize}


%%----------------------------------------------------------------------------
\subsection{Notas y comentarios}

La práctica deberá funcionar en el entorno GNU/Linux (Ubuntu) del laboratorio de la asignatura con la versión de Django que se ha usado en prácticas.

La práctica deberá funcionar desde el navegador Firefox disponible en el laboratorio de la asignatura.

%%----------------------------------------------------------------------------
\subsection{Preguntas frecuentes}
\label{sec:practica-2026-05-ltaw:preguntas}

A continuación, algunas preguntas relacionadas con el enunciado de esta práctica:

\begin{itemize}

\item ¿Puedo proponer cualquier tipo de aplicación web?

  Sí, siempre que cumpla los requisitos generales (Django, APIs externas, cobertura de contenidos del curso) y sea aprobada por los profesores en la Fase 0. Las tres ideas del enunciado son sugerencias, no obligaciones.

\item ¿Es obligatorio hablar con los profesores antes de empezar?

  Sí. La Fase 0 (propuesta inicial) es obligatoria y debe recibir el visto bueno del profesorado antes de comenzar el desarrollo. Proyectos sin propuesta aprobada no podrán obtener puntuación en la parte optativa.

\item ¿Cuántas fases intermedias hay que presentar?

  Dos: Fase 1 (prototipo funcional) y Fase 2 (versión avanzada). Las fechas concretas se acordarán con los profesores al aprobar la propuesta. Las fases deben quedar documentadas en \texttt{FASES.md}.

\item ¿Es obligatorio usar NVIDIA Build o algún LLM?

  No necesariamente. El uso de un LLM solo es obligatorio si el proyecto lo requiere por su naturaleza o si el alumno lo propone en la Fase 0. Lo que sí es obligatorio es integrar al menos una API externa real (no tiene que ser un LLM). Si el proyecto usa un LLM, NVIDIA Build es la opción recomendada, pero se permite cualquier proveedor con acceso por API.

\item ¿Dónde guardo el token de NVIDIA o de cualquier otro servicio?

  Nunca en el repositorio. Debe ir en una variable de entorno o en un fichero \verb|.env| que no se suba al control de versiones (incluido en \verb|.gitignore|).

\item ¿Cómo sé si mi proyecto cubre el 70--80\,\% de los contenidos del curso?

  En la revisión de la Fase 0, los profesores orientarán al alumno sobre esto. En general, se espera que el proyecto use: recursos web con distintos métodos HTTP, plantillas Django con herencia, modelos y formularios, autenticación, CSS/Bootstrap y al menos una API externa. El README deberá indicar explícitamente qué contenidos del curso se han aplicado.

\item ¿Puedo presentar mi proyecto a los Premios de Reutilización de Datos Abiertos?

  Sí, y se valorará como parte optativa. Si el alumno participa, deberá indicarlo en el README y aportar el enlace a la candidatura o la confirmación de participación.

\item ¿El proyecto puede ser una extensión de otra práctica de la asignatura?

  No. El proyecto final debe ser un desarrollo nuevo e independiente de las prácticas anteriores.

\item ¿La práctica es individual o en grupo?

  Es individual. Se puede debatir ideas con compañeros, pero el código y el diseño de cada entrega deben ser propios y claramente diferenciables de los de otras entregas.

\item Cuando despliego en Python Anywhere, puedo conectarme a algunos sitios pero no a otros. ¿Qué está pasando?

  Las máquinas virtuales de Python Anywhere están limitadas a los sitios en su ``lista blanca'' (\emph{whitelist}). Si la API externa que usa tu proyecto no está en esa lista, la aplicación desplegada no podrá conectarse. En ese caso, asegúrate de que la base de datos del despliegue ya incluya los datos necesarios. Más detalles en el apartado~\ref{sec:practica-2026-05-ltaw:despliegue}.

\item ¿Dónde puedo realizar el despliegue?

  En cualquier servidor accesible desde Internet durante el periodo de corrección. Python Anywhere es la opción más sencilla (instrucciones\footnote{Capítulo ``Deploy!'' de Django Girls Tutorial:\\ \url{https://tutorial.djangogirls.org/en/deploy/}}, precios\footnote{PythonAnywhere Plans and Pricing:\\ \url{https://www.pythonanywhere.com/pricing/}}). También valen plataformas como Google Cloud (instrucciones\footnote{GCP Quickstart Using a Linux VM:\\ \url{https://cloud.google.com/compute/docs/quickstart-linux}}, precios\footnote{Google Compute Engine Pricing:\\ \url{https://cloud.google.com/compute/pricing}}), un servidor propio o servicios equivalentes.

\item He desplegado en PythonAnywhere y algunos ficheros estáticos (CSS, imágenes) no se cargan. ¿Qué está pasando?

  Para servir ficheros estáticos en PythonAnywhere hay que configurarlos explícitamente en la consola web, indicando la ruta local y la URL de acceso. Añade tus ficheros estáticos en esa tabla y recarga la aplicación.

\item ¿Hay que instalar librerías adicionales?

  Depende del proyecto. Si usas NVIDIA Build con \verb|requests|, no necesitas nada fuera de lo habitual. Si usas el cliente de OpenAI, instala el paquete \verb|openai|. En cualquier caso, incluye todas las dependencias en \verb|requirements.txt|.

\end{itemize}
